% 
%
% (c) 2025 Autor, ETH Zürich
%
% !TEX root = linalg_I-II_summary.tex
% !TEX encoding = UTF-8
%

\section{Vector Spaces}

\subsection{Vector Spaces and Linear Subspaces}
\begin{definition}{Vector Space}{vect_spaces}
	A \textbf{vector space} over a field $K$ is a set $V$, endowed with two operations
	\begin{align*}
		+ : V \times V \to V,& \qquad (v_1, v_2) \mapsto v_1 + v_2,\\
		\cdot : K \times V \to V,& \qquad (\alpha, v) \mapsto \alpha \cdot v
	\end{align*}
	s.t. the following conditions (axioms) hold:
	\begin{enumerate}
		\item[(V1)] $\forall v_1, v_2, v_3 \in V: v_1 + (v_2 + v_3) = (v_1 + v_2) + v_3,$
		\item[(V2)] $\exists \text{ an element } 0 \in V : 0 + v = v \qquad \forall v \in V$,
		\item[(V3)] $\forall v \in V \; \exists v' \in V : v + v' = 0$,
		\item[(V4)] $\forall v_1, v_2 \in V : v_1 + v_2 = v_2 + v_1$,
		\item[(V5)] $\forall a,b \in K\; \forall  v \in V : a \cdot (b \cdot v) = (a\cdot b) \cdot v$,
		\item[(V6)] $\forall v \in V : 1 \cdot v = v$,
		\item[(V7)] $\forall a \in K \; \forall v_1,v_2 \in V : a \cdot (v_1 + v_2) = a\cdot v_1 + a\cdot v_2$,
		\item[(V8)] $\forall a,b \in K \; \forall v \in V : (a + b)\cdot v = a\cdot v + b \cdot v$.
	\end{enumerate}
\end{definition}

\begin{lemma}{Properties of Vector Spaces}{prop_vect_spaces}
	Let $V$ be a vector space over a field $K$. Then:
	\begin{enumerate}
		\item[(a)] The $0$ vector from (V2) is unique. To avoid confusion (e.g. with $0 \in K$) we will sometimes denote it by $0_{V}$.
		\item[(b)] The element $v'$ (corresponding to $v$) from (V3) is unique. We'll denote it by $-v$. We'll also write $w - v := w + (-v)$ for all $w \in V$,
		\item[(c)] $\forall v \in V : 0 \cdot v = 0$,
		\item[(d)] $\forall a \in K : a \cdot 0 = 0$,
		\item[(e)] $\forall v \in V : -1 \cdot v = -v$,
		\item[(f)] $\forall v \in V : -(-v) = v$,
		\item[(g)] If $a \cdot v = 0$ for some $a \in K, v \in V$, then either $a = 0$ or $v = 0$,
		\item[(h)] Associativity for $+$ holds also for some of $n$ elements. $v_1 + \hdots + v_n$ makes sense no matter how we put the brackets $v_1 + (v_2 + (\hdots + (v_{n-1} + v_n)))$ or $(v_1 + v_2) + (\hdots)$ etc.
	\end{enumerate}
\end{lemma}

\subsubsection*{Example: The Coordinate Space $K^n$}
Let $K$ be a field and $n \in \mathbb{N}$.
\[
	K^n = \{(a_1, \hdots , a_n)\; |\; a_i \in K, 1 \leq i \leq n\}.
\]
$K^n$ is a vector space over $K$ if we endow it with the following operations:
\begin{itemize}
	\item[$\bullet$] Let $v = (a_1, \hdots , a_n), w = (b_1, \hdots , b_n) \in K^n$.
	\[
		v + w = (a_1, \hdots , a_n) + (b_1, \hdots , b_n) := (a_1 + b_1, \hdots , a_n + b_n),
	\]
	\item[$\bullet$] For $a \in K$
	\[
		a\cdot v = a\cdot (a_1, \hdots , a_n) := (a\cdot a_1, \hdots , a\cdot a_n)
	\]
	\item[$\bullet$] $0 := (0, \hdots , 0) \in K^n$
\end{itemize}
With these operations $K^n$ becomes a vector space over $K$.

\begin{definition}{Linear Subspace}{lin_subsp}
	Let $V$ be a vector space over $K$. A subset $W \subseteq V$ is called a \textbf{linear subspace} of $V$ (or just subspace of $V$), if the following holds
	\begin{enumerate}
		\item[(LSS1)] $W \neq \emptyset$,
		\item[(LSS2)] $\forall w_1, w_2 \in W : w_1 + w_2 \in W$,
		\item[(LSS3)] $\forall a \in K, w \in W : a\cdot w \in W$.
	\end{enumerate}
\end{definition}

\begin{lemma}{Linear Subspace is a Vector Space}{subsp_vectsp}
	Let $V$ be a vector space over $K$ and $W \subseteq V$ a subspace. Then $W$ is a vector space on its own, when endowed with the operations from $V$.
\end{lemma}

\begin{lemma}{Criterion for a Linear Subspace}{crit_subsp}
	Let $V$ be a vector space over $K$, and $W  \subseteq V$ a subset. Then $W$ is a linear subspace of $V$ if and only if the following holds:
	\begin{enumerate}
		\item[(1)] $0_V \in W$,
		\item[(2)] $\forall a_1, a_2 \in K, w_1, w_2 \in W : a_1 w_1 + a_2 w_2 \in W$.
	\end{enumerate}
\end{lemma}

\subsubsection*{Example: Spaces of functions}
Fix a field $K$ and a non-empty set $S$. Define
\[
	K^S = \{f:S \to K \;|\; f \text{ is a function}\}.
\]
Define $+$ and $\cdot$ on $K^S$ as follows. Let $f,g \in K^S$, $a \in K$. Define
\begin{align*}
	(f+g):S \to K,& \qquad S \owns x \overset{f+g}{\longmapsto} f(x) + g(x),\\
	(a \cdot f): S \to K,& \qquad S \owns x \overset{a\cdot f}{\longmapsto} a\cdot f(x).
\end{align*}

Now Let $S = [0, 1]$ and take $K = \mathbb{R}$. Define $C(S) := \{f \in \mathbb{R}^S \;|\; f \text{ is continuous on }[0,1]\} \subseteq \mathbb{R}^S$.
Then $C(S) \subseteq \mathbb{R}^S$ is a linear subspace.

\subsection{Span of Sets}

\begin{definition}{Span}{span}
	Let $S \subseteq V$ be a subset. We define
	\[
		\operatorname{Sp}(S) := \bigcap_{W \in \mathcal{N}} W,
	\]
	where
	\[
		\mathcal{N} := \{W \;|\; W \text{ is a subspace of } V \text{ and } S \subseteq W\},
	\]
	i.e., the collection of all subspaces $W \subseteq V$ that contain $S$.
	$\operatorname{Sp}(S)$ is called the span of $S$ (in $V$).
\end{definition}

\begin{lemma}{Properties of the Span}{prop_span}
	\begin{enumerate}
		\item[(1)] $\operatorname{Sp}(S) \subseteq V$ is a linear subspace.
		\item[(2)] If $W \subseteq V$ is a linear subspace and $S \subseteq W$, then $\operatorname{Sp}(S) \subseteq W$.
	\end{enumerate}
	In other words, $\operatorname{Sp}(S)$ is a linear subspace of $V$, and moreover among those subspaces that contain $S$, $\operatorname{Sp}(S)$ is the smallest one.
\end{lemma}

\begin{definition}{Linear Combination}{lin_comb}
	Let $n \in \mathbb{N}$, $a_1, \hdots , a_n \in K$, $v_1, \hdots , v_n \in V$. A \textbf{linear combination} of the vectors $v_1, \hdots , v_n$ with coefficients $a_1, \hdots , a_n$ is the vector
	\[
		a_1 v_1 + \hdots + a_n v_n = \sum_{i=1}^{n} a_i v_i \in V.
	\]
\end{definition}
It is important to note that in this course consider only finitely many elements in the sums of linear combinations.

\begin{lemma}{Alternate Definition of Span}{alt_span}
	Let $\emptyset \neq S \subseteq V$. Then
	\[
		\operatorname{Sp}(S) = \{a_1 v_1 + \hdots + a_n v_n \;|\; n \in \mathbb{N}, a_i \in K, v_i \in S \quad \forall\, 1 \leq i \leq n\}.
	\]
\end{lemma}

\begin{remark}
	The set $S$ can be finite or infinite. Regardless of that, each linear combination uses only finite numbers of elements from $S$.
\end{remark}

Let $n \in \mathbb{N}$, $v_1, \hdots , v_n \in V$. We'll write
\[
	\operatorname{Sp}(v_1, \hdots , v_n) := \operatorname{Sp}(\{v_1, \hdots , v_n\}).
\]

\begin{remark}
	It obviously holds that
	\[
		\operatorname{Sp}(v_1, \hdots , v_n) = \bigg\{\sum_{i=1}^{n}\alpha_i v_1 \;\bigg|\; \alpha_i \in K \quad \forall \, 1\leq i \leq n\bigg\}.
	\]
\end{remark}

\begin{remark}
	It follows from the definition of span that
	\[
		\operatorname{Sp}(\emptyset) = \{0_V\}.
	\]
\end{remark}

\begin{definition}{Generating Set}{gen_set}
	Let $V$ be a vect. sp. over $K$ and $S \subseteq V$. We say that $S$ \textbf{generates} (or \textbf{spans}) $V$ if $\operatorname{Sp}(S) = V$. $S$ is called a generating set of $V$ if $\operatorname{Sp}(S) = W$, where $W \subseteq V$ is a subspace of $V$. We say that $S$ \textbf{generates} (or \textbf{spans}) $W$.
\end{definition}

\begin{definition}{Finite Dimensional}{f_dim}
	A vector space $V$ is called \textbf{finite-dimensional} if $\exists$ a \textbf{finite} set $S \subseteq V$ such that $\operatorname{Sp}(S) = V$. Otherwise $V$ is called \textbf{infinite-dimensional}.
\end{definition}

\subsection{Bases of Vector Spaces}

\subsubsection{Linear Independence}
Let $V$ be a vector space over $K$.

\begin{definition}{Linear Independence}{lin_indep}
	Let $v_1, \hdots , v_n$ be a list of $n$ vectors in $V$. We say that $v_1, \hdots , v_n$ is \textbf{linearly independent} if the only linear combination of $v_1, \hdots , v_n$ that gives $0 \in V$ is the trivial one. In other words,
	\[
		\forall a_1, \hdots , a_n \in K : a_1v_1 + \hdots + a_n v_n \quad \Longrightarrow \quad a_1 = 0, \hdots a_n = 0.
	\]
	By definition, we say that the empty list of vectors $\emptyset$ is linearly independent.
	We say that a list of vectors is linearly dependent if, the list is not linearly independent.
\end{definition}

\begin{remark}
	\begin{enumerate}
		\item[(1)] 0 can always be written as a linear combination of any finite list of vectors $0 v_1 + \hdots + 0 v_n = 0$, i.e., the trivial linear combination.
		\item[(2)] If the list $v_1, \hdots , v_n$ has two vectors that are the same, say $v_k = v_l$ for $1 \leq k < l \leq n$, then $v_1, \hdots , v_n$ is \textbf{not} linearly independent. Indeed, $v_k + (-1) v_l = 0$.
		\item[(3)] The order of the vectors in the list $v_1, \hdots , v_n$ has no effect on linear independence. So, if there are no repetitions in $v_1, \hdots , v_n$ we can talk about linear independence of the set $\{v_1, \hdots , v_n\}$.
	\end{enumerate}
\end{remark}

\begin{definition}{Linear Independence in the Context of Families}{lin_indep_fam}
	Let $\mathcal{F} = \{v_i\}_{i \in I}$ be a familiy of vectors in $V$, indexed by a set $I$. We say that $\mathcal{F}$ is linearly independent if $\forall n \in \mathbb{N}$ and any sequence of distinct indices $i_1, \hdots , i_n \in I$, the list of vectors $v_{i_1}, \hdots , v_{i_n}$ is linearly independent.
	Of course, if in $\mathcal{F}$ there are any repetitions of vectors, i.e., if $\exists i,j, i\neq j$, s.t. $v_i = v_j$, then $\mathcal{F}$ cannot be linearly independent.
\end{definition}

\begin{lemma}{Subsets of Linearly Independent Sets are Linearly Independent}{lin_indep_subset}
	If $\emptyset \neq S \subseteq V$ is linearly independent. Then every non-empty subset $S' \subseteq S$ is also linearly independent.
\end{lemma}




















